\chapter{Test curves and geodesic rays}
\label{chap:q15}

Let $X$ be a connected compact K\"ahler manifold of dimension $n\geq1$, let $\theta$ be a smooth closed real $(1,1)$-form representing a big class, and fix a normalized model potential $\phi$ of mass $m>0$. A geodesic ray records a potential at each time. Taking its infimum after subtracting a linear slope instead records a singularity type at each slope. Kiselman's minimum principle keeps these infima plurisubharmonic; the model-envelope condition recovers the maximality of the original ray. The one-variable duality behind this construction is proved below in the precise half-line form required here.

\section{Concave families of model potentials}
\label{sec:model-test-curves}
Write $M(q)=\int_X\langle(\theta+\ddc q)^n\rangle$ in the fixed background. The normalization of model potentials supplies a common upper bound, even when their singularities vary with the parameter.

\begin{definition}\label{def:model-test-curve}
A model test curve based at $\phi$ is a real number $b$ and a decreasing, pointwise concave family $\Gamma_\tau\in\PSH(X,\theta)$, $\tau<b$, such that
\[
 P_\theta[\Gamma_\tau]=\Gamma_\tau,\qquad
             (\sup_{\tau<b}\Gamma_\tau)^*=\phi.
\]
We call $b$ its terminal slope and write $\operatorname{TC}(X,\theta,\phi)$ for these curves. A curve is bounded if $\Gamma_\tau=\phi$ for all sufficiently negative $\tau$; write $\operatorname{TC}^\infty$ for this subclass. Extend a curve by
\[
 \Gamma_b=\inf_{\tau<b}\Gamma_\tau,\qquad
                \Gamma_\tau=-\infty\quad(\tau>b).
\]
The terminal value is allowed to have zero mass.
\end{definition}
\index{test curve!model}\index{test curve!bounded}\index{test curve!terminal slope}

\Needspace{10\baselineskip}
\begin{proposition}\label{prop:test-curve-parameter-continuity}
Every interior value has positive mass. The terminal value is a noncollapsed fixed point of $P_\theta$, and $\sup_X\Gamma_\tau=0$ for $\tau\leq b$. For $\tau<b$,
\[
 \Gamma_\tau=\inf_{s<\tau}\Gamma_s
                   =(\sup_{\tau<s<b}\Gamma_s)^*.
\]
The extended scalar function $\tau\mapsto\Gamma_\tau(x)$ is upper semicontinuous and concave for every $x$. The mass function decreases, tends to $m$ as $\tau\to-\infty$, and is continuous on $(-\infty,b]$ from the interior, with the corresponding one-sided interpretation at $b$.
\end{proposition}
\begin{proof}
Fixed points have supremum zero. Their decreasing terminal limit is a potential and a fixed point by \cref{prop:decreasing-model-mass-continuity}. Outside its pluripolar pole set, all the scalar functions are finite on $(-\infty,b)$ and therefore continuous there by the elementary chord inequalities for concave functions. The first displayed limit is a decreasing potential limit; the second is a canonical increasing envelope. They agree almost everywhere with $\Gamma_\tau$, hence everywhere by canonical representatives. The first equality, together with decreasing order, proves scalar upper semicontinuity at interior parameters. The definition gives it at $b$, and beyond $b$ it is immediate. Concavity survives the extension by $-\infty$.

Increasing non-pluripolar continuity gives $M(\Gamma_s)\to m$ as $s\to-\infty$. Given $\tau<b$, choose $s<\tau$ with positive mass and $r\in(\tau,b)$. Write $\tau=(1-t)r+ts$, with $0<t<1$. Concavity, mass monotonicity and positive multilinearity give
\[
 M(\Gamma_\tau)\geq M((1-t)\Gamma_r+t\Gamma_s)
                                \geq t^nM(\Gamma_s)>0.
\]
Mass monotonicity gives decreasing order. The first parameter limit has convergence of masses by decreasing model continuity; the second by increasing continuity. This proves interior continuity and the terminal assertion.
\end{proof}

\begin{example}\label{ex:constant-and-associated-test-curves}
The constant curve $\Gamma_\tau=\phi$ for $\tau\leq b$, extended by $-\infty$ afterwards, has terminal slope $b$. More generally, for $u\in\PSH(X,\theta)$ with $u\leq\phi$ define a bounded curve with terminal slope zero by
\[
 \Gamma_\tau=
 \begin{cases}
 \phi,&\tau\leq-1,\\
 P_\theta[(1+\tau)u-\tau\phi],&-1<\tau<0.
 \end{cases}
\]
The inner family is constant and then decreasing affine; since $u\leq\phi$, its chord slopes also decrease across the join at $-1$. Envelope concavity in \cref{prop:two-envelope-algebra} gives concavity of its projections. The coefficient of $\phi$ is positive at every interior parameter, so positive multilinearity and mass preservation give positive mass and the model property. The value at $-\infty$ is $\phi$. At zero its limit may have zero mass.
\end{example}
\index{test curve!associated with a potential}

\section{The half-line Legendre transform}
\label{sec:half-line-legendre-duality}
The transforms use positive times and all real slopes. Their pointwise form matters at polar points, where a spatial upper envelope can change a raw supremum.

\begin{lemma}\label{lem:half-line-scalar-legendre-duality}
Let $q:\mathbb R\to[-\infty,C]$ be upper semicontinuous, decreasing and concave, with $q=-\infty$ above a finite $b$. For $t>0$ put
\[
 g(t)=\sup_{\tau\in\mathbb R}(q(\tau)+t\tau).
\]
Then
\[
 q(\tau)=\inf_{t>0}(g(t)-t\tau).
\]
Conversely, if $g$ is finite and convex on $(0,\infty)$ and $g(t)\leq C+bt$, then, for $q(\tau)=\inf_{t>0}(g(t)-t\tau)$,
\[
 g(t)=\sup_{\tau<b}(q(\tau)+t\tau).
\]
The identically $-\infty$ case has both transforms identically $-\infty$.
\end{lemma}
\begin{proof}
For the first assertion, suppose $q$ is not identically $-\infty$. Its finite domain is an interval unbounded to the left. At an interior finite point, the chord slopes give a supporting affine upper line of slope $s\leq0$. If $s<0$, choosing $t=-s$ gives equality in the inverse formula. If $s=0$, the value is the global supremum of $q$; as $t\downarrow0$, $g(t)$ tends to that supremum. Indeed its upper bound is $\sup q+tb$ and each fixed finite value gives the reverse lower limit.

Let $a\leq b$ be the right endpoint of the finite domain. At $\tau>a$, the bound $g(t)\leq C+ta$ gives the inverse value $-\infty$ as $t\to\infty$. At $a$, upper semicontinuity gives, for each finite $L>q(a)$, an interval $(a-\delta,a]$ where $q<L$. Splitting the supremum there and on $(-\infty,a-\delta]$ gives
$g(t)-ta\leq\max\{L,C-t\delta\}$.
Letting $t\to\infty$ and then $L\downarrow q(a)$ proves equality, including $q(a)=-\infty$. The opposite inverse bound always follows from the defining supremum.

For the converse, the convex chord inequalities give a supporting affine lower line at every $t>0$. Its slope $\tau$ is at most $b$, by comparing the line with $C+bs$ as $s\to\infty$. If $\tau<b$, that line gives $q(\tau)=g(t)-t\tau$. If $\tau=b$, let $\tau_j\uparrow b$. The functions $g(s)-s\tau_j$ decrease to $g(s)-sb$, and the two infima over $s,j$ commute. Thus $q(\tau_j)\downarrow q(b)$, which again gives equality using slopes below $b$. Every defining infimum gives the reverse inequality. The infinite-value case follows directly from the formulas.
\end{proof}

\begin{definition}\label{def:test-curve-legendre-transform}
For $\Gamma\in\operatorname{TC}(X,\theta,\phi)$ set
\[
 \check\Gamma_t(x)=\sup_{\tau<b}(\Gamma_\tau(x)+t\tau),\quad t>0,
                           \qquad\check\Gamma_0=\phi.
\]
For a ray $\ell$ based at $\phi$ set
\[
 \widehat\ell_\tau(x)=\inf_{t>0}(\ell_t(x)-t\tau).
\]
The first formula is a raw supremum. Its value at zero is specified separately, since the raw initial supremum may differ from $\phi$ on a pluripolar set.
\end{definition}
\index{Legendre transform!of a test curve}\index{Legendre transform!of a geodesic ray}

\Needspace{10\baselineskip}
\begin{proposition}\label{prop:test-curve-legendre-subgeodesic}
The raw transform is jointly upper semicontinuous on $X\times(0,\infty)$ and defines a subgeodesic ray starting at $\phi$. For every $t>0$,
\[
 \check\Gamma_t\in\mathcal E(X,\theta,\phi),\qquad
 \sup_X\check\Gamma_t=bt,\qquad
 \Gamma_\tau=\inf_{t>0}(\check\Gamma_t-t\tau).
\]
The inverse formula holds for every real $\tau$. Also $\check\Gamma_t-bt$ decreases with $t$.
\end{proposition}
\begin{proof}
Write $F_t$ for the raw transform and $G$ for the canonical upper envelope of its product complexifications. Each $\Gamma_\tau+t\tau$ is a subgeodesic. The bound $F_t\leq\phi+bt$ and a fixed competitor give a noncollapsed, locally bounded-above envelope, so $G$ is $p^*\theta$-psh on the right half-plane. Imaginary-translation invariance is preserved. The raw and canonical envelopes agree almost everywhere on the product. Fubini and continuity of finite convex fibres show that $F_t=G_t$ for all positive times at almost every point of $X$. Here the terminal pole set can be removed to make $F_t$ finite at every positive time.

For $\tau<b$ put $\chi_\tau=\inf_{t>0}(G_t-t\tau)$. Kiselman's minimum principle makes it $\theta$-psh or identically $-\infty$. The scalar lemma gives $\chi_\tau\geq\Gamma_\tau$, so it is a potential; the preceding almost-everywhere equality gives $\chi_\tau=\Gamma_\tau$ almost everywhere, hence everywhere. At $b$, commuting the decreasing infima gives $\chi_b=\Gamma_b$; for $\tau>b$ both are identically $-\infty$ by $G_t\leq\phi+bt$. Applying the converse scalar lemma to each fibre of $G$ now gives $G_t=F_t$ everywhere. A finite fibre has the required convexity and upper linear bound; an identically $-\infty$ fibre is covered by the last case of the lemma. This proves the asserted upper semicontinuity and inverse formula.

We have $F_t\geq\Gamma_\tau+t\tau$ for every fixed $\tau<b$. Hence mass monotonicity gives $M(F_t)\geq M(\Gamma_\tau)$, and letting $\tau\to-\infty$ gives $M(F_t)\geq m$. The upper bound $F_t\leq\phi+bt$ gives the opposite mass bound and the relative order condition. Thus $F_t$ has relative full mass. The supremum is at most $bt$; it is at least $t\tau$ since $\sup_X\Gamma_\tau=0$. Letting $\tau\uparrow b$ proves equality. The formula
$F_t-bt=\sup_{\tau<b}(\Gamma_\tau-t(b-\tau))$
gives decreasing order.

Normalized compactness and $\sup_XF_t=bt$ give potential cluster limits as $t\downarrow0$. The upper bound puts each below $\phi$, while every fixed lower competitor puts it above $\Gamma_\tau$. Their canonical supremum is $\phi$, so the limit is $\phi$ in $L^1$. Interior continuity follows from the subgeodesic fibre argument in \cref{prop:subgeodesic-convexity-and-continuity}; the same integrable lower and upper bounds apply on every compact time interval. Thus the initial extension is a subgeodesic ray.
\end{proof}

\section{Model curves and weak geodesic rays}
\label{sec:model-test-curve-ray-correspondence}
A replacement of a ray segment changes each dual potential by a bounded amount. The model condition therefore prevents any increase of that segment.

\begin{theorem}\label{thm:model-test-curve-geodesic-ray-correspondence}
The two transforms give inverse bijections
\[
 \operatorname{TC}(X,\theta,\phi)
                  \longleftrightarrow\mathcal R(X,\theta,\phi).
\]
They restrict to a bijection between $\operatorname{TC}^\infty$ and $\mathcal R^\infty$. The terminal slope is $\sup_X\ell_1$.
\end{theorem}
\begin{proof}
\enlargethispage{3\baselineskip}
Let $F=\check\Gamma$. Replace its segment on $[a,d]$, $0\leq a<d$, by the Perron segment between its endpoints, and call the resulting subgeodesic ray $H$. Full-mass Perron existence and \cref{lem:subgeodesic-operations-and-pasting} apply, and $H\geq F$. Convexity bounds $H_t\leq bt$ throughout the replaced interval and this bound already holds elsewhere. Put $\chi_\tau=\inf_{t>0}(H_t-t\tau)$. For $\tau<b$, Kiselman's principle and $\chi_\tau\geq\Gamma_\tau$ make it a potential. Letting $t\downarrow0$ in the upper bound gives $\chi_\tau\leq0$.

Since $F_t-bt$ decreases, for $a\leq s\leq d$,
\[
 F_d-d\tau\leq F_s-s\tau+(d-a)(b-\tau).
\]
Outside that interval $H_s=F_s$, while $\chi_\tau\leq H_d-d\tau=F_d-d\tau$. Taking the infimum over all $s>0$ therefore gives
\[
 \Gamma_\tau\leq\chi_\tau
                  \leq\Gamma_\tau+(d-a)(b-\tau).
\]
The upper bound and $\chi_\tau\leq0$ make $\chi_\tau$ a competitor for the model envelope $P_\theta[\Gamma_\tau]=\Gamma_\tau$. Hence equality holds. Scalar duality gives $H=F$ everywhere. Every segment is therefore Perron, so $F$ is a weak ray.

Conversely let $\ell\in\mathcal R(X,\theta,\phi)$ and $b=\sup_X\ell_1$. \Cref{prop:relative-ray-growth-and-energy} gives $\ell_t\leq\phi+bt$ and $\sup_X\ell_t=bt$. Convexity makes $\ell_t-bt$ decrease: a convex fibre bounded above on the half-line has no positive chord slope. The decreasing family has supremum zero, so its infimum $\widehat\ell_b$ is a noncollapsed potential by \cref{cor:global-decreasing-nets}. Kiselman's principle and $\widehat\ell_\tau\geq\widehat\ell_b$ give potentials for $\tau<b$; for $\tau>b$ the upper linear bound gives the identically $-\infty$ value. The dual family is decreasing and concave. Letting $t\downarrow0$ in its defining bound shows $\widehat\ell_\tau\leq0$.

Fix $\tau<b$ and $C>0$, and put
$q=P_\theta(\widehat\ell_\tau+C,V_\theta)$.
Since $\widehat\ell_\tau\leq\ell_1-\tau\preceq\phi$, we have $q\leq\widehat\ell_\tau+C\preceq\phi$. As $q\leq0$, it is a competitor for $P_\theta[\phi]=\phi$, so $q\leq\phi$. The affine curve $q-Cs$ on $0\leq s\leq1$ lies below the endpoints of the Perron segment $\ell_{sT}-sT\tau$, for every $T>0$. Thus
\[
 q-Cs\leq\ell_{sT}-sT\tau.
\]
For fixed $s>0$, take the infimum over $T>0$ to obtain $q-Cs\leq\widehat\ell_\tau$; then let $s\downarrow0$. The reverse rooftop bound follows from $\widehat\ell_\tau\leq0$, so $q=\widehat\ell_\tau$. Letting $C\to\infty$ proves $P_\theta[\widehat\ell_\tau]=\widehat\ell_\tau$.

On a common full-measure set, the finite convex fibres of $\ell$ have initial limit $\phi$. Supporting lines at times decreasing to zero give
$\sup_{\tau<b}\widehat\ell_\tau=\phi$ there. Canonical upper envelopes give this equality everywhere after regularization. Thus $\widehat\ell$ is a model test curve based at $\phi$. Scalar duality gives the inverse identities for positive times and all slopes.

If $\Gamma_\tau=\phi$ for $\tau\leq c$, then $\phi+ct\leq\check\Gamma_t\leq\phi+bt$. Conversely \cref{prop:relative-ray-growth-and-energy} bounds a bounded-type ray below by $\phi+ct$. This gives $\widehat\ell_\tau\geq\phi$ for $\tau\leq c$; the initial limit gives equality.
\end{proof}
\index{Ross--Witt Nystr\"om correspondence}\index{geodesic ray!model test curve}

\section{Energy as an integral of model masses}
\label{sec:test-curve-energy}
A singularity curve measures how much mass is lost at each slope. Integrating that loss gives the energy per unit time of its geodesic ray. Monotonicity of the masses suffices for the integral argument.

\begin{definition}\label{def:test-curve-finite-energy}
For a model test curve of terminal slope $b$ put
\[
 \mathbf E^\phi(\Gamma)
       =bm+\int_{-\infty}^b\bigl(M(\Gamma_\tau)-m\bigr)\,\mathrm d\tau
                                      \in[-\infty,bm].
\]
The integral is the negative of a nonnegative Lebesgue integral. A test curve has finite energy if this value exceeds $-\infty$; write $\operatorname{TC}^1(X,\theta,\phi)$ for these curves. Every bounded test curve has finite energy; its initial slope is $\sup\{\tau\leq b:\Gamma_\tau=\phi\}$.
\end{definition}
\index{test curve!finite energy}\index{Monge--Amp\`ere energy!of a test curve}

\begin{lemma}\label{lem:bounded-test-curve-decreasing-approximation}
Suppose $\Gamma$ has terminal slope zero. For $\varepsilon>0$ set
\[
 \Gamma^\varepsilon_\tau=
 \begin{cases}
 \phi,&\tau\leq-\varepsilon^{-1},\\
 P_\theta[(1+\varepsilon\tau)\Gamma_\tau
                         -\varepsilon\tau\phi],&-\varepsilon^{-1}<\tau<0.
 \end{cases}
\]
These are bounded model test curves based at $\phi$. As $\varepsilon\downarrow0$, their potentials and masses decrease to those of $\Gamma$ at every interior slope, their positive-time transforms decrease pointwise to $\check\Gamma$, and
\[
 \mathbf E^\phi(\Gamma^\varepsilon)\downarrow
                                             \mathbf E^\phi(\Gamma).
\]
\end{lemma}
\begin{proof}
\enlargethispage{2\baselineskip}
On the common finite locus put $A_\tau=\phi-\Gamma_\tau$. It is nonnegative, increasing and convex. The product of two such scalar convex functions is convex: their chord upper bounds multiply, and the difference between the chord of the product and that product of chords is
$s(1-s)(A_c-A_a)(B_c-B_a)\geq0$.
Apply this to $A_\tau$ and $B_\tau=\max\{1+\varepsilon\tau,0\}$. Thus the unprojected family $\phi-B_\tau A_\tau$ is decreasing and concave, including across the join. Canonical representatives extend these inequalities everywhere. Envelope concavity preserves them. The positive coefficient of $\phi$ in each interior affine combination gives positive mass; its envelope is therefore model. The constant tail gives the required base.

The inner potentials decrease to $\Gamma_\tau$ as $\varepsilon\downarrow0$, and so do their envelopes, initially to a fixed point $q\geq\Gamma_\tau$. For a fixed interior slope, positive multilinearity shows that the masses of the inner potentials tend to $M(\Gamma_\tau)$. Mass preservation and decreasing model continuity give $M(q)=M(\Gamma_\tau)>0$. The normalized maximum characterization in \cref{thm:positive-mass-model-characterization} forces $q=\Gamma_\tau$. This proves the slope and mass limits. Monotone convergence of $m-M(\Gamma^\varepsilon_\tau)$ proves the energy limit.

Their rays decrease and dominate $\check\Gamma$, while each is at most $\phi$. The product limit is noncollapsed, has relative full-mass slices, and has initial $L^1$ limit $\phi$ by this sandwich. On every finite interval, decreasing endpoint continuity in \cref{prop:perron-geodesic-order-and-limits} makes the limit a weak ray $L$. Commuting the infima over time and $\varepsilon$ gives
\[
 \widehat L_\tau
   =\inf_{\varepsilon>0}\inf_{t>0}
                     (\check\Gamma^\varepsilon_t-t\tau)
   =\inf_{\varepsilon>0}\Gamma^\varepsilon_\tau=\Gamma_\tau
                                  \qquad(\tau<0).
\]
The inverse correspondence therefore gives $L=\check\Gamma$ pointwise.
\end{proof}

\Needspace{12\baselineskip}
\begin{theorem}\label{thm:test-curve-ray-energy-identity}
For every model test curve, including those of infinite negative energy,
\[
 E_\theta^\phi(\check\Gamma_t)=t\mathbf E^\phi(\Gamma)
                                        \qquad(t>0).
\]
Consequently the Legendre correspondence restricts to inverse bijections
\[
 \operatorname{TC}^1(X,\theta,\phi)
                   \longleftrightarrow\mathcal R^1(X,\theta,\phi),
\]
and it preserves radial energy.
\end{theorem}
\begin{proof}
Replacing $\Gamma_\tau$ by $\Gamma_{\tau+b}$ makes its terminal slope zero, subtracts $bt$ from its transform and subtracts $bm$ from its curve energy. The energy shift formula therefore reduces the proof to $b=0$.

First suppose $\Gamma_\tau=\phi$ for $\tau\leq c<0$. Fix $t>0$, $\delta=2^{-N}$ and $K=\lfloor c/\delta\rfloor$. For $K\leq j\leq0$ put
\[
 F_j=\max_{K\leq k\leq j}(\Gamma_{k\delta}+tk\delta).
\]
Every $F_j$ lies between $\phi+tK\delta$ and $\phi$, hence has bounded relative type. We claim
\begin{equation}\label{eq:test-curve-grid-energy-increment}
 t\delta M(\Gamma_{j\delta})
       \leq E(F_j)-E(F_{j-1})
       \leq t\delta M(\Gamma_{(j-1)\delta})
                                \qquad(K<j\leq0),
\end{equation}
where $E=E_\theta^\phi$.

Let $U=\{F_j>F_{j-1}\}$ off the common poles. Concavity of the sampled scalar function $\Gamma_\tau+t\tau$ shows that on $U$ its last increment is positive, and all preceding increments are at least as large. Thus on $U$,
\[
 F_j=\Gamma_{j\delta}+tj\delta,\quad
 F_{j-1}=\Gamma_{(j-1)\delta}+t(j-1)\delta,\quad
                  0<F_j-F_{j-1}\leq t\delta.
\]
The set is plurifinely open. Locality identifies its two measures with those of the corresponding model values. The measure of $\Gamma_{j\delta}$ is carried by its zero set, by \cref{prop:p-envelope-contact-concentration}; this includes the terminal value. On that set every earlier model value is zero, so it lies in $U$ and the increment is exactly $t\delta$. The lower tangent-energy bound in \cref{lem:relative-energy-tangent-and-order-bounds} therefore gives the lower bound in \eqref{eq:test-curve-grid-energy-increment}. For the upper bound, the upper tangent bound, locality on $U$ and the displayed increment bound give at most $t\delta M(\Gamma_{(j-1)\delta})$. The omitted poles have zero measure.

Since $F_K=\phi+tK\delta$, summing gives
\[
 t\delta\sum_{j=K+1}^0\bigl(M(\Gamma_{j\delta})-m\bigr)
       \leq E(F_0)
       \leq t\delta\sum_{j=K}^{-1}\bigl(M(\Gamma_{j\delta})-m\bigr).
\]
These are lower and upper sums for the bounded monotone mass function on $[K\delta,0]$. Their difference is at most $tm\delta$. That function equals $m$ below $c$, so both sums tend to $t\int_{-\infty}^0(M(\Gamma_\tau)-m)\,\mathrm d\tau$. No concavity of its logarithm is required.

The same $F_0$ is the maximum over all grid points $k\delta\leq0$, since the constant tail makes earlier values smaller. Dyadic refinement therefore increases $F_0$ to the raw transform: interior scalar continuity approximates every slope, and the terminal value contributes no larger value. Increasing energy continuity in \cref{prop:relative-energy-increasing-continuity} proves the desired identity for bounded curves.

For an arbitrary curve use the bounded approximations of the preceding lemma. Their transforms decrease to $\check\Gamma_t$, and their curve energies decrease to $\mathbf E^\phi(\Gamma)$. Extended decreasing energy continuity in \cref{prop:relative-energy-decreasing-continuity} passes the bounded identity to the limit, even when the value is $-\infty$. The finite-energy characterization and \cref{prop:relative-ray-growth-and-energy} now give the asserted restriction and radial equality.
\end{proof}

\begin{corollary}\label{cor:associated-test-curve-ray}
The bounded curve associated with $u\leq\phi$ in \cref{ex:constant-and-associated-test-curves} transforms into the singularity-associated ray $\ell^u$ of \cref{thm:relative-singularity-associated-ray}. Its energy is
\[
 \mathbf E^\phi(\Gamma)
 =\frac1{n+1}\sum_{j=0}^n
       \int_X\left\langle\theta_u^j\wedge\theta_\phi^{n-j}\right\rangle-m.
\]
\end{corollary}
\begin{proof}
Each long segment defining $\ell^u$ dominates the constant competitor $u$ and the affine competitor $\phi-t$. Their limit therefore dominates $\max\{u,\phi-t\}$. Its inverse potentials dominate $(1+\tau)u-\tau\phi$ for $-1<\tau<0$, and dominate $\phi$ for $\tau\leq-1$. This also forces its terminal slope to be zero. Those inverse potentials are model, so they dominate the projected curve $\Gamma$. Duality gives $\ell^u\geq\check\Gamma$.

Mass preservation, positive multilinearity and the elementary integral
\[
 \int_0^1 a^j(1-a)^{n-j}\,\mathrm da
                      =\frac{j!(n-j)!}{(n+1)!}
\]
give the displayed energy formula. The integral follows by $j$ integrations by parts followed by integration of $(1-a)^n$. It equals the radial energy of $\ell^u$ computed in \cref{thm:relative-singularity-associated-ray}. Both rays have bounded type; their ordered chordal distance is their radial energy difference, by \cref{thm:relative-chordal-ray-completeness}. That difference is zero, so the rays coincide.
\end{proof}

\section{Multiplier-ideal model curves}
\label{sec:i-model-test-curves}
The Legendre correspondence keeps ordinary model singularities at each slope. Requiring those models to be determined by multiplier ideals gives a smaller class, stable under the natural projection and functorial operations.

\begin{definition}\label{def:i-model-test-curve}
A model test curve is $\mathcal I$-model if
\[
 P_\theta[\Gamma_\tau]_{\mathcal I}=\Gamma_\tau
                                      \qquad(\tau<b).
\]
Write $\operatorname{TC}_{\mathcal I}(X,\theta,\phi)$ for this subclass and use superscripts $1$ and $\infty$ for its finite-energy and bounded subclasses.
Omitting $\phi$ means the union over all normalized positive-mass model bases, for either $\operatorname{TC}(X,\theta)$ or $\operatorname{TC}_{\mathcal I}(X,\theta)$.
\end{definition}
\index{test curve!I-model@$\mathcal I$-model}

\begin{proposition}\label{prop:i-model-test-curve-boundaries}
The base and terminal potential of an $\mathcal I$-model test curve are $\mathcal I$-model. In particular this subclass is empty unless $\phi$ is $\mathcal I$-model.
\end{proposition}
\begin{proof}
Increasing multiplier-envelope continuity in \cref{prop:increasing-two-envelope-continuity} gives the base assertion. At the terminal slope, monotonicity gives
\[
 \Gamma_b\leq P_\theta[\Gamma_b]_{\mathcal I}
       \leq\inf_{\tau<b}P_\theta[\Gamma_\tau]_{\mathcal I}
       =\inf_{\tau<b}\Gamma_\tau=\Gamma_b.
\]
The terminal argument also applies when its mass is zero.
\end{proof}

\begin{proposition}\label{prop:test-curve-background-and-ideal-projections}
Let $\omega$ be a smooth closed semipositive $(1,1)$-form, allowing zero. Applying either envelope levelwise gives model curves
\[
 \tau\longmapsto P_{\theta+\omega}[\Gamma_\tau],\qquad
 \tau\longmapsto P_{\theta+\omega}[\Gamma_\tau]_{\mathcal I},
\]
with the same terminal slope and respective bases
$P_{\theta+\omega}[\phi]$ and $P_{\theta+\omega}[\phi]_{\mathcal I}$.
The second curve is $\mathcal I$-model. In one background its projection is monotone and idempotent, dominates the original curve, and fixes every $\mathcal I$-model curve. Both constructions preserve the bounded subclass.
\end{proposition}
\begin{proof}
The enlarged class is big. Positive multilinearity gives every original interior potential positive mass in it. Envelope monotonicity and concavity preserve decreasingness and concavity of the parameter family. Mass preservation for the ordinary envelope, and mass monotonicity for the multiplier envelope, give positive interior mass and the model property. Increasing envelope continuity identifies the bases. Multiplier idempotence gives the second curve's $\mathcal I$-model property. Pointwise order, idempotence and fixed points follow from the corresponding potential statements in \cref{prop:multiplier-envelope-projection,prop:two-envelope-algebra}. A constant initial tail remains constant after projection.
\end{proof}

\begin{proposition}\label{prop:test-curve-smooth-background-change}
If $\theta'=\theta+\ddc g$ for smooth real $g$, set
\[
 \phi'=P_{\theta'}[\phi-g],\qquad
                 \Gamma'_\tau=P_{\theta'}[\Gamma_\tau-g].
\]
This is a bijection from $\operatorname{TC}(X,\theta,\phi)$ to $\operatorname{TC}(X,\theta',\phi')$. It preserves terminal slopes, masses, curve energies, boundedness and finite energy, and restricts to the $\mathcal I$-model subclasses. Successive changes compose.
\end{proposition}
\begin{proof}
Smooth subtraction identifies the currents. The bounds in \cref{prop:envelopes-change-background-modification} give
\[
 \Gamma_\tau-g+\inf_Xg\leq\Gamma'_\tau
                              \leq\Gamma_\tau-g+\sup_Xg.
\]
Thus each projected value has the bounded type and mass of $\Gamma_\tau-g$. Envelope algebra preserves parameter concavity and order; increasing continuity gives base $\phi'$. Applying the reverse construction gives a normalized model of the bounded type of $\Gamma_\tau$, hence $\Gamma_\tau$ itself. The same argument proves composition for two smooth changes. Mass equality at every slope and at the base preserves the energy integral; constant tails preserve boundedness.

An interior $\mathcal I$-model value has positive mass and is good. Smooth subtraction preserves goodness. Its ordinary projection in the new background therefore equals its multiplier projection by \cref{prop:i-good-background-independence}. This proves the restricted bijection.
\end{proof}

\begin{proposition}\label{prop:test-curve-modification-invariance}
For a proper bimeromorphic morphism $\pi:Y\to X$ of connected compact K\"ahler manifolds, pointwise pullback gives a bijection
\[
 \operatorname{TC}(X,\theta,\phi)
       \longleftrightarrow
              \operatorname{TC}(Y,\pi^*\theta,\pi^*\phi).
\]
It preserves terminal slopes, masses, curve energies, boundedness and finite energy, and restricts to the $\mathcal I$-model subclasses.
\end{proposition}
\begin{proof}
The pulled-back class is big. The potential bijection of \cref{prop:smooth-target-psh-modification} preserves order and concavity in both directions, since $\pi$ is surjective. Both envelopes commute with pullback by \cref{prop:envelopes-change-background-modification}. Their countable selections and $L^1$-continuous pullback also identify the canonical initial suprema. Top-degree mass invariance in \cref{prop:nonpluripolar-modification-pushforward} preserves every interior mass and the base mass. The energy formula then proves energy invariance; constant tails and multiplier projection give the other assertions.
\end{proof}

\begin{definition}\label{def:i-model-geodesic-ray}
A weak geodesic ray is $\mathcal I$-model if its inverse Legendre curve is $\mathcal I$-model. The correspondence identifies these rays with $\operatorname{TC}_{\mathcal I}$, also in the finite-energy and bounded subclasses. This condition specifies the multiplier-ideal data of every slope.
\end{definition}
\index{geodesic ray!I-model@$\mathcal I$-model}

\section{Curves from section filtrations}
\label{sec:section-filtration-test-curves}
A multiplicative filtration gives a singularity curve by taking normalized section weights. The section asymptotic formula identifies its negative-slope tail and supplies its positive interior masses.

\begin{definition}\label{def:multiplicative-section-filtration}
Let $L$ have a smooth quadratic Hermitian metric $h$ with curvature $\theta$, and assume $\{\theta\}$ is big. Its section ring is
$R(X,L)=\bigoplus_{k\geq0}H^0(X,L^k)$, with the product of sections. A bounded multiplicative filtration consists of subspaces $\mathcal F_k^\lambda\subseteq H^0(X,L^k)$, decreasing and left-continuous in $\lambda$, such that
\[
 \mathcal F_k^\lambda\mathcal F_l^\eta
                       \subseteq\mathcal F_{k+l}^{\lambda+\eta},\qquad
 \mathcal F_k^{-Ck}=H^0(X,L^k),\quad
 \mathcal F_k^{Ck}=0\quad(k\geq1)
\]
for a fixed $C>0$. Left-continuity means intersection over smaller parameters. In degree zero set $\mathcal F_0^\lambda=\mathbb C$ for $\lambda\leq0$ and zero otherwise. Put
\[
 a_k=\sup\{\lambda:\mathcal F_k^\lambda\ne0\},\qquad
              b=\sup_{k\geq1}\frac{a_k}{k},
\]
with $a_k=-\infty$ when $H^0(X,L^k)=0$.
\end{definition}
\index{filtration!of a section ring}\index{filtration!critical slope}

\begin{lemma}\label{lem:filtration-critical-slope}
The critical slope $b$ is finite and equals $\lim_{k\to\infty}a_k/k$. For every $\tau<b$, at least one $\mathcal F_k^{k\tau}$ is nonzero.
\end{lemma}
\begin{proof}
The big-class unweighted section formula in \cref{cor:big-or-zero-line-bundle-volume} gives nonzero sections for every sufficiently large $k$, since the big class has positive volume. For those degrees $-Ck\leq a_k\leq Ck$. The top parameter is attained: along parameters increasing to a finite $a_k$, the nonzero nested subspaces stabilize in finite dimension, and left-continuity keeps their intersection nonzero. Products of nonzero sections are nonzero, so $a_{k+l}\geq a_k+a_l$ whenever both terms are finite.

Choose $k_0$ after which all terms are finite. Fix any $h$ with $a_h$ finite and write $k=qh+r$ with $k_0\leq r<k_0+h$. Superadditivity gives $a_k\geq qa_h+a_r$, hence
$\liminf_k a_k/k\geq a_h/h$.
Take the supremum over such $h$; the reverse upper bound is the definition of $b$. This proves the limit and finiteness. If $\tau<b$, some $a_k/k>\tau$, and the definition gives the last assertion.
\end{proof}

\Needspace{14\baselineskip}
\begin{theorem}\label{thm:filtration-model-test-curve}
For $\tau<b$ define
\[
 \rho_\tau=
 \left(\sup_{\substack{k\geq1,\ s\in\mathcal F_k^{k\tau}\\
                     \sup_X|s|_{h^k}^2\leq1}}
                      \frac1k\log|s|_{h^k}^2\right)^*.
\]
Then
\[
 \Gamma_\tau=P_\theta[\rho_\tau]
                 =P_\theta[\rho_\tau]_{\mathcal I}
\]
is a bounded $\mathcal I$-model test curve based at $V_\theta$, with terminal slope $b$. Its ray is bounded and $\mathcal I$-model. The construction is unchanged by replacing $h$ by another smooth metric on $L$, under the preceding smooth-background bijection.
\end{theorem}
\begin{proof}
A nonzero filtered section at some degree supplies a noncollapsed lower competitor. Every normalized section weight is $\theta$-psh by Poincar\'e--Lelong and at most zero, so $\rho_\tau$ is a potential. The family decreases with $\tau$.

For rational $s=p/q\in(0,1)$, take sections $s_i\in\mathcal F_{k_i}^{k_i\tau_i}$ with supremum norm at most one. Their product
\[
 s_0^{k_1(q-p)}s_1^{k_0p}
       \in\mathcal F_{k_0k_1q}^{\,k_0k_1q((1-s)\tau_0+s\tau_1)}
\]
has norm at most one, and its normalized logarithm is the corresponding convex combination of the two section weights. Taking suprema proves concavity for rational weights. For $\tau_0<\tau_1$, weights decreasing to a given real weight, together with decreasingness in $\tau$, prove the same inequality for all real weights. Canonical representatives give concavity everywhere.

At a fixed nonzero degree, the section supremum has rational analytic singularities. Indeed a basis identifies its supremum-norm unit ball with a bounded set containing a Euclidean ball; its weight differs by a bounded function from $k^{-1}\log\sum|s_i|_{h^k}^2$. Thus each degree weight is good, and their canonical supremum $\rho_\tau$ is good by \cref{prop:i-good-positive-algebra-and-suprema}. Initially define $\Gamma_\tau=P_\theta[\rho_\tau]_{\mathcal I}$. Multiplier projection makes it a normalized fixed point of both envelopes, even before its positive mass is known. Envelope algebra preserves concavity and decreasingness.

For $\tau<-C$ all filtered spaces are full, so $\rho_\tau$ and $\Gamma_\tau$ are constant. Write $q$ for this common projected value. For every section of $L^k$, normalization gives
$kq\geq\log|s|_{h^k}^2$ up to a constant depending on the section. Hence $|s|^2\mathrm e^{-kq}$ is locally bounded almost everywhere and
\[
 H^0(X,L^k\otimes\mathcal I(kq))=H^0(X,L^k).
\]
The big-class weighted and unweighted section formulas therefore give
\[
 M(q)=\operatorname{vol}_{\mathcal I}(\theta_q)
              =\operatorname{vol}(\{\theta\})=M(V_\theta)>0.
\]
Here $q$ is its own multiplier envelope, and the last equality is \cref{lem:minimal-boundary-volume-attainment}. As $q\leq V_\theta$ is a normalized ordinary fixed point of this mass, the model maximum characterization gives $q=V_\theta$.

The concave model family now has this positive-mass initial tail. The interior-positivity argument of \cref{prop:test-curve-parameter-continuity} gives every interior value positive mass. Goodness of $\rho_\tau$ implies its mass equals that of its multiplier envelope in this background, by \cref{prop:i-good-background-independence}. Positive-mass envelope equality therefore gives
$P_\theta[\rho_\tau]=P_\theta[\rho_\tau]_{\mathcal I}$,
proving the stated ordinary formula. The bounded correspondence gives the ray assertion.

Finally changing the smooth metric changes each section weight by the same smooth function $-g$. The squared supremum norms have ratios between $\mathrm e^{-k\|g\|_\infty}$ and $\mathrm e^{k\|g\|_\infty}$. Thus the section envelopes differ by a bounded function after subtraction of $g$. The background-envelope bounds put $P_{\theta'}[\rho_\tau-g]$ in the bounded type of $\Gamma_\tau-g$. Its normalized model projection is therefore $P_{\theta'}[\Gamma_\tau-g]$, proving metric independence.
\end{proof}

\begin{example}\label{ex:critical-filtration-degree-oscillation}
On $\mathbb P^1$ with $L=\mathcal O(1)$, give the full degree-$k$ space weight $a_k=0$ for even $k$ and $a_k=-1$ for odd $k$: set $\mathcal F_k^\lambda=H^0(X,L^k)$ for $\lambda\leq a_k$ and zero otherwise. The inequalities $a_{k+l}\geq a_k+a_l$ make this a bounded multiplicative filtration, with critical slope zero. At slope zero its even-degree spaces are full and its odd-degree spaces vanish. For the Fubini--Study metric, a normalized power of a linear form peaks at any prescribed point; Cauchy--Schwarz bounds its norm by one. The degree section suprema at the critical slope are therefore zero in even degrees and $-\infty$ in odd degrees. Their full-degree sequence has no pointwise limit. At every negative slope the curve of the theorem is zero, and its canonical terminal value is zero as well.
\end{example}
\index{filtration!critical-degree oscillation}
